%\begin{sloppypar}
%\hyphenation{RANDRIAMAMPANDRY}
\indent \justifying Tout d'abord, je tiens à remercier M. Rivoarison RANDRIANASOLO, Professeur Titulaire, Président de l'Université d'Antananarivo ; M. Tendro RADANIELINA, Maître de conférences, Doyen de la Faculté des Sciences ; et M. Dimbimalala RANDRIANASOLOHARISOA, Maître de conférences, Chef du Département de Physique et Applications, de m'avoir permis de soutenir ce mémoire de master.\\

%\indent Mes sincères remerciements vont également à M. Minoson RAKOTOMALALA, Professeur Titulaire, d'avoir accepté de présider le jury. Je remercie aussi M. Solohery RANDRIAMAMPANDRY, Professeur Associ\'e, ainsi que M.Toky RANDRIAMAMPANDRY, Maître de conférences, d'avoir accepté de faire partie du jury en tant qu'examinateurs.\\
\indent Mes sincères remerciements vont également à M. ~Minoson RAKOTOMALALA, professeur titulaire, pour avoir accepté d'être le président du jury. Je remercie aussi M.~Solohery \linebreak RANDRIAMAMPANDRY, Professeur, ainsi que M. ~Toky RANDRIAMAMPANDRY, Maître de conférences, d'avoir fait partie du jury en tant qu'examinateurs. Je tiens aussi à les remercier pour leur enseignement et leur soutien pendant la période de mes études à l'université.\\

\indent Je suis également reconnaissant envers mon encadreur, M. Roger IANJAMASIMANANA, Chercheur Postdoctoral, pour le temps qu'il a consacré à ce projet. Sans son aide, ses remarques et ses suggestions, je n'aurais pas pu mener ce mémoire à terme.\\

\indent Enfin, j'exprime ma gratitude à mes parents pour leur soutien constant tout au long de mes années d'études, ainsi que durant la recherche et la rédaction de ce mémoire. Cet accomplissement n'aurait pas été possible sans eux. Je remercie également mes amis et collègues pour leurs encouragements.

%\end{sloppypar}